% Context Formers
The rules for context formation are given as follows:
    \begin{center}
        \begin{bprooftree}
            \AxiomC{$\phantom{\Gamma, x : A \vdash M : B}$}
            \UnaryInfC{$\vdash \diamond \; \Ctxt$}
        \end{bprooftree}\quad
        \begin{bprooftree}
            \AxiomC{$\vdash \Gamma \; \Ctxt$}
            \AxiomC{$\Gamma \vdash A \Type$}
            \BinaryInfC{$\vdash \Gamma, x : A \; \Ctxt$}
        \end{bprooftree}
    \end{center}
    \begin{center}
        \begin{bprooftree}
            \AxiomC{$\vdash \Gamma, x : A, \Delta \; \Ctxt$}
            \UnaryInfC{$\Gamma, x:A , \Delta \vdash x:A$}
        \end{bprooftree}\quad
        \begin{bprooftree}
            \AxiomC{$\Gamma, x : A\vdash b(x):B(x)$}
            \AxiomC{$\Gamma\vdash a:A$}
            \BinaryInfC{$\Gamma\vdash b(a):B(a)$}
        \end{bprooftree}
    \end{center}
\vspace{10pt}

The rules for the empty-, unit-, and boolean type are as follows, where their elimination and computation principles are left implicit.
% empty, unit, bool
\begin{center}
    \begin{bprooftree}
        \AxiomC{$\phantom{\Gamma \vdash M : \Pi x : A.\,B}$}
        \UnaryInfC{$\Gamma \vdash \bot \Type$}
    \end{bprooftree}\quad
    \begin{bprooftree}
        \AxiomC{$\phantom{\Gamma \vdash M : \Pi x : A.\,B}$}
        \UnaryInfC{$\Gamma \vdash \top \Type$}
    \end{bprooftree}\quad
    \begin{bprooftree}
        \AxiomC{$\phantom{\Gamma \vdash M : \Pi x : A.\,B}$}
        \UnaryInfC{$\Gamma \vdash \mathsf{Bool} \Type$}
    \end{bprooftree}\quad
\end{center}
\begin{center}
    \begin{bprooftree}
        \AxiomC{\phantom{$\Gamma, x : A \vdash b : B$}}
        \UnaryInfC{$\Gamma \vdash \star : \top$}
    \end{bprooftree}\quad
    \begin{bprooftree}
        \AxiomC{\phantom{$\Gamma, x : A \vdash b : B$}}
        \UnaryInfC{$\Gamma \vdash \mathsf{tt} : \mathsf{Bool}$}
    \end{bprooftree}\quad
    \begin{bprooftree}
        \AxiomC{\phantom{$\Gamma, x : A \vdash b : B$}}
        \UnaryInfC{$\Gamma \vdash \mathsf{ff} : \mathsf{Bool}$}
    \end{bprooftree}\quad
\end{center}
\vspace{10pt}

The rules for $\Pi$-types include definitional $\beta$- and $\eta$-rules.
% Pi type
\begin{center}
    \begin{bprooftree}
        \AxiomC{$\Gamma \vdash A \Type$}
        \AxiomC{$\Gamma, x : A \vdash B(x) \Type$}
        \BinaryInfC{$\Gamma \vdash \Pi x : A.\,B(x) \Type$}
    \end{bprooftree}\quad
\end{center}
\begin{center}
    \begin{bprooftree}
        \AxiomC{$\Gamma, x : A \vdash b(x) : B(x)$}
        \UnaryInfC{$\Gamma \vdash \lambda x \: A \. b(x): \Pi x : A.\,B(x)$}
    \end{bprooftree}\quad
    \begin{bprooftree}
        \AxiomC{$\Gamma \vdash f : \Pi x : A.\,B(x)$}
        \AxiomC{$\Gamma \vdash a : A$}
        \BinaryInfC{$\Gamma \vdash f \,a : B(a)$}
    \end{bprooftree}
\end{center}
\begin{center}
    \begin{bprooftree}
        \AxiomC{$\Gamma, x : A \vdash b(x) : B(x)$}
        \AxiomC{$\Gamma \vdash a : A$}
        \BinaryInfC{$\Gamma \vdash (\lambda x \: A \. b(x)) \, a \equiv b(a): B(a)$}
    \end{bprooftree}\quad
    \begin{bprooftree}
        \AxiomC{$\Gamma \vdash f : \Pi x : A.\,B(x)$}
        \UnaryInfC{$\Gamma \vdash (\lambda x \: A \. f \, x) \equiv f : \Pi x : A.\,B(x)$}
    \end{bprooftree}
\end{center}
\vspace{10pt}

Similarly to $\Pi$-types, $\Sigma$-types also come with $\beta$- and $\eta$-rules. The $\eta$-rule for $\Sigma$-types is defined via its destructors, where $\pi_1$ and $\pi_2$ are the projections derived from the elimination principle. 
\begin{center}
    \begin{bprooftree}
        \AxiomC{$\Gamma \vdash A \Type$}
        \AxiomC{$\Gamma, x : A \vdash B(x) \Type$}
        \BinaryInfC{$\Gamma \vdash \Sigma x:A.\,B(x) \Type$}
    \end{bprooftree}\quad
    \begin{bprooftree}
        \AxiomC{$\Gamma \vdash a : A$}
        \AxiomC{$\Gamma \vdash b : B(a)$}
        \BinaryInfC{$\Gamma \vdash \< a, b \> : \Sigma x:A.\,B(x)$}
    \end{bprooftree}
    \end{center}
    \begin{center}
        \begin{bprooftree}
            \AxiomC{$\Gamma, z : \Sigma x:A.\,B(x) \vdash P(z)\Type$}
            \noLine
            \UnaryInfC{$\Gamma, x : A, y : B \vdash p(x,y) : P(\< x, y \>)$}
            \UnaryInfC{$\Gamma, z : \Sigma x:A.\,B(x) \vdash \ind_\Sigma(z, p(x,y)) : P(z)$}
        \end{bprooftree}
    \end{center}
    \begin{center}
        \begin{bprooftree}
            \AxiomC{$\Gamma, z : \Sigma x:A.\,B(x) \vdash P(z)\Type$}
            \noLine
            \UnaryInfC{$\Gamma, x : A, y : B \vdash p : P(\< x, y \>)$}
            \UnaryInfC{$\Gamma, x : A, y : B \vdash \ind_\Sigma(\< x, y \>, p) \equiv p(x,y) : P(\< x, y \>)$}
        \end{bprooftree}
    \end{center}
    \begin{center}
        \begin{bprooftree}
            \AxiomC{$\Gamma \vdash c : \Sigma x:A.\,B(x)$}
            \UnaryInfC{$\Gamma \vdash c \equiv \< \pi_1(c), \pi_2(c) \> : \Sigma x:A.\,B(x)$}
        \end{bprooftree}
    \end{center}
\vspace{10pt}

Propositional identity types are defined with the standard induction principle.
% Identity Type
    \begin{center}
        \begin{bprooftree}
            \AxiomC{$\Gamma \vdash A \Type$}
            \AxiomC{$\Gamma \vdash a_1, a_2 : A$}
            \BinaryInfC{$\Gamma \vdash a_1 =_A a_2 \Type$}
        \end{bprooftree}\quad
        \begin{bprooftree}
            \AxiomC{$\Gamma \vdash A \Type$}
            \AxiomC{$\Gamma \vdash a : A$}
            \BinaryInfC{$\Gamma \vdash \refl \, a : a =_A a$}
        \end{bprooftree}\quad
    \end{center}
    \begin{center}
        \begin{bprooftree}
            \AxiomC{$\Gamma , x : A, y: A, z : x =_A y \vdash P(x, y, z) \Type$}
            \noLine
            \UnaryInfC{$\Gamma , x : A \vdash p(x) : P(x, x, \refl \, x)$}
            \UnaryInfC{$\Gamma, x:A, y:A, z: x=_A y \vdash \ind_=(x,y, z, p) : P(x, y, z)$}
        \end{bprooftree}
    \end{center}
    \begin{center}
        \begin{bprooftree}
            \AxiomC{$\Gamma , x : A, y : A, z : x =_A y \vdash P(x, y, z) \Type$}
            \noLine
            \UnaryInfC{$\Gamma , x : A \vdash p(x) : P(x, x, \refl \, x)$}
            \UnaryInfC{$\Gamma, x:A \vdash \ind_=(x, x, \refl \, x, p) \equiv p(x) : P(x, x, \refl \, x)$}
        \end{bprooftree}
    \end{center}
\vspace{10pt}

The universe type is closed under the base types, propositional equality, and the type constructors. The definitional type equalities for the small types ensure that they lift to the proper type.
    \begin{center}
        \begin{bprooftree}
        \AxiomC{\phantom{$\Gamma \vdash u : \U$}}
        \UnaryInfC{$\Gamma \vdash \U \Type$}
        \end{bprooftree}\quad
        \begin{bprooftree}
        \AxiomC{$\Gamma \vdash A : \U$}
        \UnaryInfC{$\Gamma \vdash \El A \, \Type$}
        \end{bprooftree}
    \end{center}
    \begin{center}
        \begin{bprooftree}
            \AxiomC{$\phantom{\Gamma \vdash}$}
            \UnaryInfC{$\Gamma \vdash \bot^\U : \U$}
        \end{bprooftree}\quad
        \begin{bprooftree}
            \AxiomC{$\phantom{\Gamma \vdash}$}
            \UnaryInfC{$\Gamma \vdash \top^\U : \U$}
        \end{bprooftree}\quad
        \begin{bprooftree}
            \AxiomC{$\phantom{\Gamma \vdash}$}
            \UnaryInfC{$\Gamma \vdash \mathsf{Bool}^\U : \U$}
        \end{bprooftree}
    \end{center}
    \begin{center}
        \begin{bprooftree}
            \AxiomC{$\Gamma \vdash A : U$}
            \AxiomC{$\Gamma, x : A \vdash B(x) : \U$}
            \BinaryInfC{$\Gamma \vdash \Pi^U x:A.\, B(x): \U$}
        \end{bprooftree}
    \end{center}
    \begin{center}
        \begin{bprooftree}
            \AxiomC{$\Gamma \vdash A : \U$}
            \AxiomC{$\Gamma, x : \El A \vdash B(x) : \U$}
            \BinaryInfC{$\Gamma \vdash \Sigma^\U x:A.\, B(x) : \U$}
        \end{bprooftree}
    \end{center}
    \begin{center}
        \begin{bprooftree}
            \AxiomC{$\Gamma \vdash A : \U$}
            \AxiomC{$\Gamma\vdash a : \El A$}
            \AxiomC{$\Gamma\vdash a' : \El A$}
            \TrinaryInfC{$\Gamma \vdash a=^U_Aa' : \U$}
        \end{bprooftree}
    \end{center}
%Type eqs
    \begin{center}
    \begin{align*}
        \El (\bot^\U) &\equiv \bot & \El (\Sigma^U x:A.\, B(x)) &\equiv \Sigma_{x:\El A} (\El B(x)) \\
        \El (\top^\U) &\equiv \top & \El (\Pi^U x:A.\, B(x)) &\equiv \Pi_{x:\El A} (\El B(x)) \\
        \El (\mathsf{Bool}^\U) &\equiv \mathsf{Bool} & \El(a=^U_Aa')&\equiv(a=_{\El A}a')
    \end{align*}
    \end{center}
\vspace{10pt}

The type of sizes is defined with constructors for zero and successor. The ordering on sizes is defined as a mere proposition, with constructors for reflexivity and transitivity.
% Sizes
\begin{center}
    \begin{bprooftree}
        \AxiomC{\phantom{$\Gamma \vdash s_1, s_2 : \Size$}}
        \UnaryInfC{$\Gamma \vdash \Size \Type$}
    \end{bprooftree}\quad
    \begin{bprooftree}
        \AxiomC{$\Gamma \vdash  i, j : \Size$}
        \UnaryInfC{$\Gamma \vdash \; i \leq j : U$}
    \end{bprooftree}\quad
    \begin{bprooftree}
        \AxiomC{\phantom{$\Gamma \vdash s_1, s_2 : \Size$}}
        \UnaryInfC{$\Gamma \vdash 0 : \Size$}
    \end{bprooftree}\quad
    \begin{bprooftree}
        \AxiomC{$\Gamma \vdash i  : \Size$}
        \UnaryInfC{$\Gamma \vdash {\uparrow \! i} : \Size$}
    \end{bprooftree}
\end{center}

\begin{center}
    \begin{bprooftree}
        \AxiomC{$\Gamma\vdash i:\Size$}
        \AxiomC{$\Gamma\vdash j:\Size$}
        %\RightLabel{\(\le\)-form,}
        \BinaryInfC{$\Gamma\vdash i\le j\Type$}
    \end{bprooftree}
    \begin{bprooftree}
        \AxiomC{$\Gamma \vdash i : \Size$}
        %\RightLabel{\(\le\)-intro\(_0\),}
        \UnaryInfC{$\Gamma \vdash{\sf{le}_0}\,i:\El(0 \le i)$}
    \end{bprooftree}\quad
    \begin{bprooftree}
        \AxiomC{$\Gamma \vdash i : \Size$}
        %\RightLabel{\(\le\)-intro\(_{\sf{suc}}\),}
        \UnaryInfC{$\Gamma \vdash{\sf{le}_{\sf{suc}}}\,i:\El(i \le {\uparrow}i) $}
    \end{bprooftree} \\[2ex]
    \begin{bprooftree}
        \AxiomC{$\Gamma \vdash i : \Size$}
        %\RightLabel{\(\le\)-intro\(_{\sf{refl}}\),}
        \UnaryInfC{$\Gamma \vdash\sf{le}_{\sf{refl}}\,i:\El(i \le i)$}
    \end{bprooftree}\quad
    \begin{bprooftree}
        \AxiomC{$\Gamma \vdash p:i\le j$}
        \AxiomC{$\Gamma \vdash q:j\le k$}
        %\RightLabel{\(\le\)-intro\(_{\sf{trans}}\).}
        \BinaryInfC{$\Gamma \vdash\sf{le}_{\sf{trans}}\,p\,q:\El(i \le k)$}
    \end{bprooftree} \\[2ex]
    \begin{bprooftree}
        \AxiomC{$\Gamma \vdash p:i\le j$}
        \AxiomC{$\Gamma \vdash q:i\le j$}
        \BinaryInfC{$\Gamma \vdash\sf{le}_{\sf{eq}}\,p\,q :\El(p=q)$}
    \end{bprooftree}
\end{center}
\vspace{10pt}

The type of the existential quantifier over sizes is defined as follows. Note that the elimination principle is limited to small types.
    \begin{center}
        \begin{bprooftree}
            \AxiomC{$\Gamma, i : \Size \vdash A(i) : \U$}
            \UnaryInfC{$\Gamma \vdash \exists i . A(i) : \U $}
        \end{bprooftree}
    \end{center}
    \begin{center}
        \begin{bprooftree}
            \AxiomC{$\Gamma, i : \Size \vdash A(i) : \U$}
            \AxiomC{$\Gamma \vdash s : \Size$}
            \AxiomC{$\Gamma \vdash a : \El A(i)$}
            \TrinaryInfC{$\Gamma \vdash \langle s, a \rangle^\exists : \El (\exists i . A(i))$}
        \end{bprooftree}
    \end{center}
    \begin{center}
        \begin{bprooftree}
            \AxiomC{$\Gamma, z : \El (\exists i. A(i)) \vdash P(z) : \U$}
            \noLine
            \UnaryInfC{$\Gamma, i : \Size, x : \El A \vdash p(i,x) : \El P(\langle i, x \rangle^\exists)$}
            \UnaryInfC{$\Gamma, z : \El(\exists i.A(i))\vdash \ind_\exists(z,p(i,x)) : \El P(z)$}
        \end{bprooftree}
    \end{center}
    \begin{center}
        \begin{bprooftree}
            \AxiomC{$\Gamma, z : \El (\exists i. A(i)) \vdash P(z) : \U$}
            \noLine
            \UnaryInfC{$\Gamma, i : \Size, x : \El A \vdash p : \El P(\langle i, x \rangle^\exists)$}
            \noLine
            \UnaryInfC{$\Gamma\vdash s:\Size$\qquad$\Gamma\vdash a:\El A(s)$}
            \UnaryInfC{$\Gamma\vdash \ind_\exists(\langle s, a \rangle^\exists, p(i,x)) \equiv p(s,a) : \El P(\langle s,a \rangle^\exists)$}
        \end{bprooftree}
    \end{center}
\vspace{10pt}

The universal quantifier is defined analogously to the $\Pi$-type, only as an impredicative quantifier.
\begin{center}
    \begin{bprooftree}
        \AxiomC{$\Gamma, i : \Size \vdash A(i) :U$}
        \UnaryInfC{$\Gamma \vdash \forall i . A(i) :U$}
    \end{bprooftree}\quad
\end{center}
\begin{center}
    \begin{bprooftree}
        \AxiomC{$\Gamma, i : \Size \vdash a(i) : A(i)$}
        \UnaryInfC{$\Gamma \vdash \lambda^\forall i : \Size \. a(i): \El (\forall i . A(i)) $}
    \end{bprooftree}\quad
    \begin{bprooftree}
        \AxiomC{$\Gamma \vdash f : \El(\forall i . A(i))$}
        \AxiomC{$\Gamma \vdash s : \Size$}
        \BinaryInfC{$\Gamma \vdash f \, s : \El A(s)$}
    \end{bprooftree}
\end{center}
\begin{center}
    \begin{bprooftree}
        \AxiomC{$\Gamma, i : \Size \vdash a(i) : \El A(i)$}
        \AxiomC{$\Gamma \vdash s : \Size$}
        \BinaryInfC{$\Gamma \vdash (\lambda^\forall i \: \Size \. a(i)) \, s \equiv a(s): \El A(s)$}
    \end{bprooftree}\quad
    \begin{bprooftree}
        \AxiomC{$\Gamma \vdash f : \El(\forall i . A(i))$}
        \UnaryInfC{$\Gamma \vdash (\lambda^\forall i \: \Size \. f \, i) \equiv f : \El (\forall i . A(i))$}
    \end{bprooftree}
\end{center}
\vspace{10pt}

Finally, the fixpoint operator on sizes, corresponding to well-founded induction, is defined as follows. Note that the associated $\beta$-rule is propositional.
    \begin{center}
        \begin{bprooftree}
            \AxiomC{$\Gamma, i : \Size \vdash A(i) \Type$}
            \AxiomC{$\Gamma \vdash f : \forall i . ((\forall j < i. A(j)) \to A(i))$}
            \BinaryInfC{$\Gamma \vdash \fix f : \forall i. A(i)$}
        \end{bprooftree}
    \end{center}
    \begin{center}
    \begin{bprooftree}
            \AxiomC{$f : \forall i . ((\forall j < i. A(j)) \to A(i))$}
            \UnaryInfC{$\Gamma \vdash \fix^\beta \, f : \forall i \. \fix f \, i = f \, i \, (\lambda j < i\. \fix f  \, j)$}
        \end{bprooftree}
    \end{center}